\section{Introduction}\label{sec:introduction}

This paper asks a bounded foundational question: once a base structure is fixed, can a genuinely \emph{internal} operator create new admissible operator power on that same base? Here ``internal'' means carrier-preserving and not essentially dependent on witnesses, selectors, evaluators, or ambient objects imported from outside the original regime. The question arises because many standard methods---completion, saturation, forcing, universal constructions, and canonicalization---produce conclusions that are later read back to an original subject matter. The issue is not whether those methods are legitimate or fruitful. It is whether, relative to a fixed base, they amount to new internal operator power.

The central claim is negative and deliberately narrow. It is best read as a fixed-base exhaustion result rather than as a global anti-generative thesis. On a fixed base structure equipped with an admissibility regime, every purported carrier-preserving strengthening is either conservative on old-base-relevant content or else depends essentially on external data and so is not internal to that base. In particular, a carrier-preserving extension cannot totalize an operation that was genuinely partial on old tuples. The result is therefore a boundary theorem about internality, not a global thesis that external mathematics is dispensable.

The argument has three parts. First, in a partial-structure setting the paper proves a rigidity theorem: if an old primitive becomes newly defined on an old tuple, then the induced old structure has changed, so the move was not carrier-preserving after all. Second, it proves an auxiliary-data dichotomy: once the effect of a purported operator is mediated by a selector, witness, evaluator, completion boundary, or comparison map, one asks whether that datum is already fixed by the base. If it is, the operator is conservative; if not, the operator is external. Third, it proves an exhaustion theorem for the main kinds of apparent internal strengthening and then applies that scheme to several standard families.

The intended contrast is thus between internal generation and external method. Forcing works by adjoining generic objects and passing to a richer ambient universe; metric completion passes to a completed object; saturation and compactness work via larger models or witness structures; universal constructions are posed relative to ambient comparison data \citep{Hodges1993,Marker2002,Jech2003,Kanamori2008,MacLane1998}. None of this is challenged here. The claim is only that, measured against a fixed original base, such methods are external or scope-extending unless the final readback is conservative on the old base.

A terminological caution is useful. ``Admissibility'' is used here in a local, regime-relative sense, not as a synonym for admissibility in the sense of admissible set theory \citep{Barwise1975}. The comparison notion used throughout is simply conservativeness on old-base-relevant content: an operator or redescription adds no new values for old operations on old tuples and no new consequences about the old base beyond those already fixed by the original regime. In particular cases that notion may be sharpened to a more specific standard relation, such as definitional extension, conservative extension over the old signature, or interdefinability.

\paragraph*{Informal theorem.}
If a purported operator on a fixed base preserves the carrier and depends only on data already fixed by that base, then it is conservative on old-base-relevant content. If it requires genuinely external data, it is not internal to the base. In particular, no carrier-preserving extension of a nondegenerate partial base totalizes a genuinely partial old primitive on old tuples.

\paragraph*{Scope.}
The main text develops five flagship families: selection and canonical representatives, completion, model-theoretic saturation and existence, forcing and absoluteness, and existence by universal property. In each case the task is the same: isolate the load-bearing datum or ambient move and classify it as internally fixed, conservative readback, or explicit scope change. The paper does not claim that these families are mathematically shallow. It claims that, relative to a fixed base, their generative status is exhausted by that classification.

\paragraph*{Roadmap.}
\Cref{sec:prior-work} positions the paper relative to conservative extension, theory comparison, partial structures, model theory, forcing, universal constructions, and foundational methodology. \Cref{sec:formal-setting} gives the formal framework. \Cref{sec:rigidity} proves the no-internal-totalization theorem, and \Cref{sec:auxiliary-data} proves the auxiliary-data dichotomy. \Cref{sec:exhaustion} gives the exhaustion theorem for apparent strengthening mechanisms. \Cref{sec:case-studies} develops the flagship cases, \Cref{sec:no-emergence} shows that recombination does not generate a new family, and \Cref{sec:main-theorem} states the final synthesis. Appendix~\ref{app:extended-cases} records longer case analyses and mixed-family examples.
